\section{Model}
\label{sec:model}

\policyengine{} UK runs on \texttt{policyengine\_core}, a fork of
OpenFisca-Core, which evaluates a directed graph of parameterized formulas
over person, benefit unit, and household entities. \todo{One paragraph on the
engine's evaluation model (entities, periods, formula/parameter separation),
citing policyengine\_core documentation or the sibling policyengine-core
paper if one exists by the time this section is drafted.}

\subsection{Program coverage}
\label{sec:model-coverage}

Table~\ref{tab:program-coverage} lists the programs \policyengine{} UK
models, drawn from the model's own coverage registry,
\texttt{policyengine\_uk/programs.yaml}, which is the single source of truth
served by the \texttt{/uk/metadata} API and rendered on the public model
coverage page. As of this writing (2026-07-04), the registry lists
\textbf{45} programs: \textbf{38} at \texttt{status: complete} and
\textbf{7} at \texttt{status: partial}, spanning \dwp{} (17, all complete:
including Universal Credit, Pension Credit, Housing Benefit, and Personal
Independence Payment), \hmrc{} (12, 11 complete and 1 partial: including
income tax, National Insurance, Child Benefit, Capital Gains Tax, Stamp
Duty Land Tax, VAT, fuel duty, and student loan repayments), the Department
for Education (9, 3 complete and 6 partial), plus single-agency coverage
for Revenue Scotland (2, complete), and one program each (all complete) for
DCMS, Local (council tax/business rates administration), Treasury, Ofgem,
and Social Security Scotland.

\begin{table}[H]
    \centering
    \caption{Program coverage by agency (source:
    \texttt{policyengine\_uk/programs.yaml}, verified 2026-07-04)}
    \label{tab:program-coverage}
    \tablefont
    \begin{tabular}{lccc}
        \toprule
        Agency & Programs & Complete & Partial \\
        \midrule
        \dwp{} & 17 & 17 & 0 \\
        \hmrc{} & 12 & 11 & 1 \\
        DfE & 9 & 3 & 6 \\
        Revenue Scotland & 2 & 2 & 0 \\
        DCMS & 1 & 1 & 0 \\
        Local & 1 & 1 & 0 \\
        Treasury & 1 & 1 & 0 \\
        Ofgem & 1 & 1 & 0 \\
        Social Security Scotland & 1 & 1 & 0 \\
        \midrule
        Total & 45 & 38 & 7 \\
        \bottomrule
    \end{tabular}
    \tablenote{Counts of programs per agency and status verified directly
    from \texttt{policyengine\_uk/programs.yaml} by parsing every
    \texttt{agency}/\texttt{status} pair (45 total entries; agency
    subtotals cross-checked against independent grep counts of
    \texttt{status: complete} (38) and \texttt{status: partial} (7)).
    Agency string is exactly \texttt{Treasury} in the source file, not
    ``HM Treasury''.}
\end{table}

Of the 7 \texttt{status: partial} programs, 1 is at \hmrc{}
(\texttt{business\_rates}, whose note reads: ``Free hours for
disadvantaged 2-year-olds and, since April 2024, working families with
younger children'' -- \todo{this note reads as though it was copied from a
different program's entry; re-check against the live file before quoting
it in prose, since it does not obviously describe business rates}) and 6
are at DfE, all English student-support programs modelled as ``first-pass''
implementations per their \texttt{notes} fields: the Childcare Grant,
Parents' Learning Allowance, Adult Dependants' Grant, Travel Grant,
Disabled Students' Allowance, and the 16 to 19 Bursary Fund. Each note
specifies a concrete simplifying assumption -- e.g. the Childcare Grant note
states it uses ``explicit full-time, SEN, and childcare-support inputs plus
maintenance-loan proxies''; the Travel Grant note states the survey does
not observe reimbursable travel costs, so the model requires explicit
placement and eligible-expense inputs instead. \todo{Quote or summarize the
remaining partial-program notes (Parents' Learning Allowance, Adult
Dependants' Grant, Disabled Students' Allowance, 16 to 19 Bursary Fund) in
full once this subsection is drafted in prose -- verified to exist in the
source file as of 2026-07-04 but not all reproduced here.}

\subsection{Devolved and regional variation}

\policyengine{} UK models devolved tax and benefit variation where the
underlying legislation devolves it: Scottish income tax bands (via
\texttt{income\_tax}'s \texttt{parameter\_prefix: gov.hmrc.income\_tax},
noted as including ``Scottish rates via devolved tax bands''), Land and
Buildings Transaction Tax (Revenue Scotland) and Land Transaction Tax (Welsh
Revenue Authority) in place of Stamp Duty Land Tax, and Scottish-specific
benefits (Scottish Child Payment, Scottish Carer's Allowance Supplement,
Scottish Child Winter Heating Assistance) that also appear as distinct line
items in the \ukmod{} comparison in Section~\ref{sec:validation-ukmod}.
\todo{Confirm Northern Ireland coverage explicitly -- programs.yaml
coverage fields say ``UK'', ``Great Britain'', or ``England and Northern
Ireland'' per program; state which programs exclude Northern Ireland and
why.}

\subsection{Known coverage gaps}

\todo{Enumerate known gaps as of v1: e.g. any programs not yet in
programs.yaml at all (rather than partial), indirect tax effects of reforms
(tracked as an open documentation item, policyengine-uk\#1114/\#1119), and
any other gaps identified by the model coverage page or open GitHub issues
tagged as scope limitations. This subsection should be exhaustive and
specific -- it is the section reviewers and users will check most closely.}
